% STATUS: written 2026-11-06 (English gate clean; sources read in full as recorded below).
% Sources read in full: Dedekind--Weber 1882 (J. reine angew. Math. 92, 181--290; scan archive.org journalfrdierei42crelgoog, Sections 11 and 23 and the plan on p. 183 checked in the OCR text; Section 23, headed "Die Differentialquotienten", opens on p. 259), Hilbert 1897 (Jahresber. DMV 4, Section 3 checked against the page image, p. 179, where the Differente is defined), Hecke 1923 (Section 36, OCR text).
% Not yet read: Hensel--Landsberg 1902; Kronecker 1882 (Grundzuege, same volume 92, pp. 1--122).
\section[Differentiation in number and function fields: Dedekind, Weber, Hensel, Hilbert][Number and function fields]{Differentiation in number and function fields: Dedekind, Weber, Hensel, Hilbert}\label{sec:origins:function-fields}

The derivative of a polynomial did algebraic work long before derivations
were defined on rings. It did this work in two places. The first place is
the theory of algebraic functions of one variable. The second place is the
arithmetic of number fields.

Riemann defined differentials on a Riemann surface by analysis. Dedekind and
Weber wanted the same objects without any use of continuity. In number
fields, the derivative of a defining polynomial gives a number called the
\emph{different}~\cite[\S 3, p.~179]{Hilbert1897}. It is defined through the derivative.

Dedekind and Weber gave an answer in 1882. Their memoir is \emph{Theorie der algebraischen Functionen einer
Ver\"anderlichen}~\cite{DedekindWeber1882}. It is written in German. It
builds the theory of Riemann surfaces again, from the arithmetic of ideals. The
introduction states the plan: \translated{A formal definition of the
differential quotient based on this then leads to the genus and to a
completely general, elegant representation of the differentials of the first
kind}{Eine darauf gegr\"undete formale Definition des Differentialquotienten
f\"uhrt sodann zu der Geschlechtszahl und zu einer ganz allgemeinen,
eleganten Darstellung der Differentiale erster Gattung}{German}%
~\cite[p.~183]{DedekindWeber1882}. The same page adds a remark. Up to the
Riemann--Roch theorem, continuity and power series play no role.

The definition is in Section~23, \emph{Die Differentialquotienten}
\cite[\S 23, p.~259]{DedekindWeber1882}. We
follow it in our translation. Let $\Omega$ be the field of algebraic
functions. Let $\alpha$ and $\beta$ be elements of $\Omega$. Assume that
$\beta$ is not constant. The two elements satisfy an irreducible equation
$F(\alpha,\beta) = 0$. Dedekind and Weber show that $\Omega$ contains a
function $(d\alpha/d\beta)$. At almost every point, its value is the limit of
the difference quotient. A footnote gives a purely algebraic formula:
\[
  \Bigl(\frac{d\alpha}{d\beta}\Bigr) = -\,\frac{F'(\beta)}{F'(\alpha)}.
\]
Here $F'(\alpha)$ and $F'(\beta)$ are the two partial derivatives of~$F$.
With this formula they prove the chain rule:
\[
  \frac{d\alpha}{d\beta} = \frac{d\alpha}{d\gamma}\,\frac{d\gamma}{d\beta}
\]
\cite[\S 23, pp.~259--260]{DedekindWeber1882}. Next they give every element
$\alpha$ a differential $d\alpha$. It satisfies $d\alpha =
(d\alpha/d\beta)\, d\beta$. They note that \translated{the differentials of the
constants, and only these, are to be set equal to zero}{Die Differentiale der
Constanten, und nur diese sind Null zu setzen}{German}. Finally, take any
rational relation $F(\alpha,\beta,\gamma,\ldots) = 0$. From it they derive
$F'(\alpha)\,d\alpha + F'(\beta)\,d\beta + F'(\gamma)\,d\gamma + \cdots = 0$%
~\cite[\S 23, p.~261]{DedekindWeber1882}.

In modern language, the construction reads as follows. Fix a nonconstant element $z$. Then the map $\alpha \mapsto d\alpha/dz$ is a
derivation of $\Omega$. Its kernel is exactly the field of constants. The map
$\alpha \mapsto d\alpha$ is an early form of the universal derivation
(Section~\ref{sec:commutative:kahler}). Dedekind and Weber did not use these
words. They needed a derivative defined by algebra alone. With it, they
reached the genus and the differentials of the first kind without analysis.

On the side of arithmetic, the derivative appears in the \emph{Zahlbericht} of Hilbert~\cite{Hilbert1897}. It is written in German.
Let $\alpha$ be an algebraic number. Let $\alpha', \ldots, \alpha^{(m-1)}$ be
its conjugates. Hilbert calls the product
\[
  \delta(\alpha) = (\alpha-\alpha')(\alpha-\alpha'')\cdots(\alpha-\alpha^{(m-1)})
\]
the \emph{different} of~$\alpha$. He notes that it is again a number of the
field. Indeed, put $f(x) = (x-\alpha)(x-\alpha')\cdots(x-\alpha^{(m-1)})$.
Then
\[
  \delta(\alpha) = \Bigl[\frac{df(x)}{dx}\Bigr]_{x=\alpha}
\]
\cite[\S 3, p.~179]{Hilbert1897}. So the different is the derivative of the
polynomial at $\alpha$. The discriminant of $\alpha$ is the product of the
squared differences of the conjugates~\cite[\S 3, p.~179]{Hilbert1897}.

Twenty-six years later, Hecke links the different of a number to the different of the field. He does
this in his lectures~\cite{Hecke1923}, which are written in German. We follow
them in our translation. He calls the ideal of Satz~101
\translated{the different or the fundamental ideal of the field}{die
Differente oder das Grundideal des K\"orpers}{German} (Satz~103). Let
$\vartheta$ be an integral generator of the field. Let $F$ be its defining
polynomial. Then $F'(\vartheta)$ is the different of $\vartheta$. Hecke
proves the following. Let $\alpha$ be a number such that $\alpha$ times the
different is integral. Then $\alpha = \Phi(\vartheta)/F'(\vartheta)$ for
some $\Phi(\vartheta)$ in the ring generated by
$\vartheta$~\cite[\S 36, pp.~133--134]{Hecke1923}.

Looking back, the two lines have much in common. Both use the derivative as an operation of polynomial algebra. No limit
is involved. Both lines also tie the derivative to the branching of an
extension. Dedekind and Weber call the ideal of multiple prime factors the
\translated{branching ideal}{Verzweigungsideal}{German}%
~\cite[\S 11, p.~226]{DedekindWeber1882}. They relate it to the differential
quotients in Section~23. The different of Hilbert is the derivative of the
defining polynomial. The next step is to treat the derivative as an abstract
operation on a field, with no polynomial behind it. This is the subject of
Section~\ref{sec:origins:fields-and-separability}.
